% Lab report.
% Units: \SI{9.81}{\metre\per\second\squared} or \qty{9.81}{m/s^2} (siunitx).
% Plots: pgfplots reads numbers from data.csv, so the figure updates with the data.
\documentclass[11pt,a4paper]{article}

\usepackage[a4paper,margin=2.5cm]{geometry}
\usepackage[english]{babel}
\usepackage{microtype}
\usepackage{amsmath}
\usepackage{siunitx}   % numbers and units
\usepackage{booktabs}
\usepackage{graphicx}
\usepackage{pgfplots}
\pgfplotsset{compat=1.18}

\usepackage[style=numeric,backend=biber]{biblatex}
\addbibresource{references.bib}
\usepackage{csquotes}
\usepackage[hidelinks]{hyperref}
\usepackage[capitalise]{cleveref}

% ---- Report details: change these ------------------------------------------
\title{Measuring the Acceleration Due to Gravity with a Simple Pendulum}
\author{Your Name \and Lab Partner\\[0.3em]\small Course, Group 3}
\date{Experiment performed on 10 March 2026}

\begin{document}
\maketitle

\begin{abstract}
  We measured the period of a simple pendulum for five lengths and obtained
  $g = \qty{9.79(6)}{m/s^2}$, in agreement with the accepted value.
\end{abstract}

\section{Aim}
To determine the local gravitational acceleration $g$ from the period of a
simple pendulum.

\section{Theory}
For small amplitudes, the period $T$ of a pendulum of length $L$ is
\cite{halliday2013}
\begin{equation}
  T = 2\pi\sqrt{\frac{L}{g}}
  \quad\Longrightarrow\quad
  T^2 = \frac{4\pi^2}{g}\,L .
  \label{eq:period}
\end{equation}
A plot of $T^2$ against $L$ is therefore a straight line with slope $4\pi^2/g$.

\section{Apparatus and method}
\begin{itemize}
  \item A steel bob of mass \qty{50}{g} on a light string.
  \item A metre rule (resolution \qty{1}{mm}) and a stopwatch (\qty{0.01}{s}).
\end{itemize}
For each length we timed 20 oscillations three times and took the mean.

\section{Results}
\Cref{tab:data} lists the measurements; \cref{fig:plot} shows $T^2$ against $L$.

\begin{table}[htbp]
  \centering
  \caption{Pendulum length and mean period.}
  \label{tab:data}
  \begin{tabular}{S[table-format=1.2] S[table-format=1.3] S[table-format=1.3]}
    \toprule
    {$L$ (\unit{m})} & {$T$ (\unit{s})} & {$T^2$ (\unit{s^2})} \\
    \midrule
    0.20 & 0.898 & 0.806 \\
    0.40 & 1.270 & 1.613 \\
    0.60 & 1.555 & 2.418 \\
    0.80 & 1.796 & 3.226 \\
    1.00 & 2.008 & 4.032 \\
    \bottomrule
  \end{tabular}
\end{table}

\begin{figure}[htbp]
  \centering
  \begin{tikzpicture}
    \begin{axis}[
        width=0.75\linewidth, height=6cm,
        xlabel={$L$ (\unit{m})}, ylabel={$T^2$ (\unit{s^2})},
        grid=major, legend pos=north west]
      \addplot[only marks,mark=*,blue] table[col sep=comma,x=L,y=T2] {data.csv};
      \addlegendentry{measured}
      \addplot[red,thick,domain=0.15:1.05] {4.033*x};
      \addlegendentry{fit: $T^2 = 4.033\,L$}
    \end{axis}
  \end{tikzpicture}
  \caption{$T^2$ against $L$ with a least-squares line through the origin.}
  \label{fig:plot}
\end{figure}

\section{Analysis}
From the slope $k = \qty{4.033(25)}{s^2/m}$ and \cref{eq:period},
\[
  g = \frac{4\pi^2}{k} = \qty{9.79(6)}{m/s^2}.
\]

\section{Discussion}
Compare with the accepted value, name the largest source of uncertainty
(here, reaction time when starting and stopping the watch) and suggest how to
reduce it.

\section{Conclusion}
State the result with its uncertainty in one sentence.

\printbibliography

\end{document}
